\uuid{2BEj}
\titre{Ensembles}
\auteur{bodin ; croizat ; sacre}
\organisation{exo7}
\licence{CC BY-NC-SA 4.0}
\datecreate{2022-03-08}
\niveau{L1}
\module{Algèbre}
\chapitre{Logique, ensemble, raisonnement}
\sousChapitre{Ensemble}
\difficulte{3}
\theme{Ensembles}
\exo7qcmid{hilisit.logique.037}

\contenu{
\begin{question}
\qclassification{100.02}
Soit $E$ un ensemble. Pour deux parties $A$ et $B$ de $E$, on définit $\Delta(A,B)=(A \cup B)\setminus (A \cap B)$. Quelles sont les affirmations vraies (quel que soit le choix de $A$ et $B$) ?
\begin{answers}
	\good{Si $A=B$, $\Delta(A,B) = \varnothing$.}
	\bad{$A\cup B\subset \Delta(A,B)$}
	\bad{$A\cap B\subset \Delta(A,B)$}
    \good{$\Delta(A,B)=(A\setminus B)\cup (B\setminus A)$}
\end{answers}
\begin{explanations}
$\Delta(A,B)$ comprend les éléments qui sont soit dans $A$ soit dans $B$, mais pas dans les deux à la fois.
Faites des schémas avec des "patates" (diagrammes de Venn) pour visualiser ces propriétés.
\end{explanations}
\end{question}
}
